\section*{Chapter \ref{ch:b3:income-and-wellbeing-evidence} --- Income, Wellbeing and What Economics Measures}

\begin{solution}{exo:b3:income-and-wellbeing-evidence:1}
Bulgaria: 4.8 at 22{,}954 dollars in 2013, 6.2 at 34{,}221 in 2024, an arrow up and to the right. Denmark: 8.0 at 60{,}746, 7.6
at 71{,}035, an arrow to the right and slightly down. % figdata: sat_table
\end{solution}

\begin{solution}{exo:b3:income-and-wellbeing-evidence:2}
$0.69\ln 1.1 = 0.066$ points at both incomes. One more dollar adds $0.69/20{,}000 = 3.5 \times 10^{-5}$ points at 20{,}000 and $0.69/80{,}000 = 8.6 \times 10^{-6}$ at 80{,}000, four
times less.
\end{solution}

\begin{solution}{exo:b3:income-and-wellbeing-evidence:3}
$(78 - 81)/81 \times 5 = -0.185$; $e^{-0.185} = 0.831$, a loss of 16.9~\% of consumption.
\end{solution}

\begin{solution}{exo:b3:income-and-wellbeing-evidence:4}
A panel following the same people after a large, lasting rise in income --- lottery
winners, say --- in which satisfaction rises and stays up for years, with no drift
back towards its earlier level.
\end{solution}

\begin{solution}{exo:b3:income-and-wellbeing-evidence:5}
The tested solution is \code{ex05_ranks.py}, in the chapter's \code{exercises/}
folder. Cyprus, Lebanon, Bosnia and Herzegovina, Jordan and Albania rise most (Cyprus from 25th to 6th); Qatar,
the United Arab Emirates, Botswana, China and Gabon fall most (Qatar from 2nd to 75th). The fallers have low consumption
shares, long hours per person (Qatar and the Emirates, with many migrant workers) or short lives (Botswana).
The risers owe much to few hours worked per person, which the measure counts as leisure even where it is
unemployment or few women in paid work: a weakness of the measure, not a finding. % output: ex05_ranks
\end{solution}

\begin{solution}{exo:b3:income-and-wellbeing-evidence:6}
$\ln\lambda = \ln(36/40) - 40(0.12^2 - 0.15^2)/2 = -0.105 + 0.162 = 0.057$: B's equivalent consumption is $1.058 \times 40{,}000 = 42{,}330$, more than A's despite lower
consumption.
\end{solution}

\begin{solution}{exo:b3:income-and-wellbeing-evidence:7}
$\sigma = \sqrt2\,\Phi^{-1}(0.70) = 0.742$ and $\sqrt2\,\Phi^{-1}(0.65) = 0.545$; the term is $(0.550 - 0.297)/2 = 0.127$, so welfare rises by $e^{0.127} - 1 = 13.5~\%$ of consumption.
\end{solution}

\begin{solution}{exo:b3:income-and-wellbeing-evidence:8}
Slope $= [(-1)(-0.9) + 0 + (1)(0.8)]/2 = 0.85$; doubling income adds $0.85\ln 2 = 0.59$ points.
\end{solution}

\begin{solution}{exo:b3:income-and-wellbeing-evidence:9}
Vignettes: the same described life rated by people in both countries. If one country
rates the same vignettes lower, part of its lower satisfaction is scale use. Movers
who keep their old ratings in a new country point the same way.
\end{solution}

\begin{solution}{exo:b3:income-and-wellbeing-evidence:10}
As in the proof of the proposition. For the Gini index, the Lorenz curve of a lognormal distribution
gives $G = 2\Phi(\sigma/\sqrt2) - 1$; solving, $\sigma = \sqrt2\,\Phi^{-1}((1 + G)/2)$.
\end{solution}

\begin{solution}{exo:b3:income-and-wellbeing-evidence:11}
The tested solution is \code{ex11_kappa.py}, in the chapter's \code{exercises/}
folder. With a year worth 3, 5 and 7: France 93.7, 103.8 and 115.0~\% of the United States; Germany 95.1, 101.4, 108.0;
Japan 60.4, 69.6, 80.1; Brazil 17.3, 16.0, 14.8; China 9.9, 9.7, 9.5; India 9.6, 7.9, 6.4. Countries whose people live longer than
in the United States gain as a year is valued more, those whose people live less lose. % output: ex11_kappa
\end{solution}

\begin{solution}{exo:b3:income-and-wellbeing-evidence:12}
For example: living standards and the augmented Solow model, which predicts
convergence only among economies with the same saving and population growth;
unconditional convergence among all countries would not refute it, but income gaps
unrelated to saving, schooling and population growth would. Prices and money and the
quantity theory, which predicts a slope of one between money growth and inflation
over long periods; a slope far from one among high-inflation countries would refute it.
\end{solution}

\begin{solution}{pb:b3:income-and-wellbeing-evidence:brenland}
% figdata: brenland_stats
\textbf{1.} 12{,}377 and 31{,}470 dalens of 2015; $\ln(31{,}470/12{,}377)/50 = 1.87~\%$ a year.

\textbf{2.} $12{,}377 \times 0.5768 = 7{,}139$ and $31{,}470 \times 0.5722 = 18{,}006$ dalens.

\textbf{3.} 1.85~\% a year.

\textbf{4.} They ignore longer lives and shorter hours, which raise welfare, and
rising inequality, which lowers it.

\textbf{5.} $\ln\lambda = (e_{75} - e_{25})u_{75}/e_{25} + \ln(c_{75}/c_{25}) - \theta(h_{75}^2 - h_{25}^2)/2 - (\sigma_{75}^2 - \sigma_{25}^2)/2$.

\textbf{6.} $\theta = 0.6 \times (1/0.5722)/0.1301^2 = 61.9$.

\textbf{7.} $900/5{,}840 = 0.154$ and $760/5{,}840 = 0.130$.

\textbf{8.} $-61.9 \times (0.154^2 - 0.130^2)/2 = -0.211$.

\textbf{9.} $\ln(7{,}139/18{,}006) = -0.925$.

\textbf{10.} $u_{75} = 5 - 0.925 - 0.211 + 0.061 = 3.925$; the term is $(72 - 81)/81 \times 3.925 = -0.436$.

\textbf{11.} The other terms compare two flows; this one compares how many years
there are, so it needs to know how much a year is worth.

\textbf{12.} 0.329.

\textbf{13.} $\sqrt2\,\Phi^{-1}(0.635) = 0.488$ and $\sqrt2\,\Phi^{-1}(0.664) = 0.600$.

\textbf{14.} $-(0.488^2 - 0.600^2)/2 = +0.061$: 1975 was more equal, which counts for 1975's welfare.

\textbf{15.} $\ln\lambda = -0.436 - 0.925 - 0.211 + 0.061 = -1.511$; $\lambda = 0.221$.

\textbf{16.} $1.511/50 = 3.02~\%$ a year.

\textbf{17.} Longer lives: 0.87 points a year; shorter hours add 0.42; rising inequality subtracts 0.12.

\textbf{18.} 2.58~\% and 3.47~\% a year.

\textbf{19.} $27{,}700/31{,}470 = 88.0~\%$.

\textbf{20.} Life $-0.119$; consumption $\ln(16{,}620/18{,}006) = -0.080$; leisure $-0.086$; inequality $-0.002$.

\textbf{21.} $e^{-0.287} = 75.0~\%$.

\textbf{22.} Life expectancy: two years less, worth $-0.119$, more than the consumption gap.

\textbf{23.} They would show whether Kestrians report better or worse lives; they
could not say whether a gap comes from income, from scale use or from things the
measure leaves out.

\textbf{24.} Brenland's \omterm{def:b3:income-and-wellbeing-evidence:welfare}{consumption-equivalent welfare} grew 3.0~\% a year from 1975 to 2025, against 1.9~\% for its GDP per person.
\end{solution}
